% !TEX root = ../proyect.tex

\chapter{Introducci\'on}\label{cap:introduccion}

\section{Contexto y Motivaci\'on}\label{sec:contexto}

La digitalizaci\'on de los peque\~nos negocios locales se ha convertido en una necesidad en la econom\'ia actual. Los salones de belleza, barber\'ias y peluquer\'ias operan mayoritariamente bajo modelos de gesti\'on tradicionales: agenda f\'isica, llamadas telef\'onicas y la memoria del profesional. Sin embargo, los clientes actuales esperan poder reservar en cualquier momento sin realizar llamadas, recibir recordatorios autom\'aticos y gestionar sus citas desde dispositivos m\'oviles.

Este proyecto nace de una necesidad real: la peluquer\'ia familiar de los autores carec\'ia de herramientas digitales adaptadas a su flujo de trabajo. Las soluciones comerciales existentes presentan dos problemas fundamentales: un coste recurrente elevado (entre 200 y 400 euros anuales en suscripciones) y una complejidad excesiva con funcionalidades innecesarias.

La asignatura de Complementos de Bases de Datos ofrece el marco id\'oneo para abordar este reto, ya que el n\'ucleo del sistema reside en el dise\~no e implementaci\'on de una base de datos relacional que soporte eficientemente todas las operaciones del negocio.

\section{Descripci\'on del Problema}\label{sec:problema}

La gesti\'on tradicional de citas en peque\~nos negocios locales presenta las siguientes deficiencias:

\textbf{P\'erdida de tiempo productivo:} Las interrupciones por llamadas telef\'onicas o mensajes de WhatsApp para agendar, modificar o cancelar citas suponen una distracci\'on continua. El profesional debe elegir entre ignorar la llamada o interrumpir el servicio que est\'a prestando, afectando negativamente a la experiencia del cliente.

\textbf{Descontrol del historial:} Sin un sistema centralizado, resulta dif\'icil recordar qu\'e servicio se le realiz\'o a un cliente hace meses. Informaci\'on cr\'itica como f\'ormulas de color, preferencias de corte o alergias se pierde o queda dispersa en notas f\'isicas.

\textbf{Falta de m\'etricas:} La ausencia de datos estructurados impide obtener una visi\'on clara del rendimiento del negocio. Preguntas como ``\textquestiondown cu\'ales son los servicios m\'as demandados?'' carecen de respuesta inmediata.

\textbf{Errores humanos:} La gesti\'on manual es propensa a solapamientos de citas (\textit{overbooking}) o huecos vac\'ios innecesarios, con impacto econ\'omico directo.

\section{Objetivos del Proyecto}\label{sec:objetivos}

El objetivo general es desarrollar una plataforma web integral que digitalice la gesti\'on de una barber\'ia o peluquer\'ia. Los objetivos espec\'ificos son:

\begin{enumerate}
    \item Dise\~nar e implementar una base de datos relacional que modele clientes, profesionales, servicios, citas y configuraciones del negocio.
    
    \item Desarrollar un portal p\'ublico de reservas accesible desde dispositivos m\'oviles que permita consultar servicios, visualizar disponibilidad y realizar reservas 24/7.
    
    \item Desarrollar un panel de administraci\'on con calendario visual, fichas de cliente con historial t\'ecnico, y gesti\'on de horarios y servicios.
    
    \item Garantizar la integridad de los datos mediante restricciones, triggers y procedimientos almacenados que eviten solapamientos de citas.
    
    \item Implementar una arquitectura \textit{single-tenant} donde el negocio sea propietario absoluto de sus datos.
\end{enumerate}

\section{Alcance del Proyecto}\label{sec:alcance}

El sistema se estructura en dos m\'odulos:

\textbf{Portal de Reservas (Cliente):} Aplicaci\'on web responsive donde los clientes pueden explorar el cat\'alogo de servicios, ver disponibilidad en tiempo real, realizar reservas, y gestionar sus citas desde un \'area personal.

\textbf{Panel de Administraci\'on (Negocio):} Herramienta privada con calendario de citas, sistema de fichas de cliente para notas t\'ecnicas, gesti\'on de horarios y vacaciones, y administraci\'on del cat\'alogo de servicios.

\textbf{Funcionalidades excluidas:} Sistema de pagos online, marketing automatizado, gesti\'on de inventario y aplicaci\'on m\'ovil nativa.

\section{Tecnolog\'ias Utilizadas}\label{sec:tecnologias}

Para el desarrollo del proyecto se han seleccionado las siguientes tecnolog\'ias:

\begin{itemize}
    \item \textbf{Base de datos:} PostgreSQL, por su robustez y soporte avanzado de restricciones e integridad referencial.
    \item \textbf{Backend:} Node.js con Express, proporcionando una API REST para la comunicaci\'on con el frontend.
    \item \textbf{Frontend:} React con TypeScript, para interfaces de usuario modernas y mantenibles.
    \item \textbf{Estilos:} Tailwind CSS, para un dise\~no responsive y consistente.
\end{itemize}

\section{Estructura de la Memoria}\label{sec:estructura}

La memoria se organiza de la siguiente forma:

\begin{itemize}
    \item \textbf{Cap\'itulo 2:} Dise\~no de la Base de Datos --- Modelo conceptual, l\'ogico y f\'isico.
    \item \textbf{Cap\'itulo 3:} Implementaci\'on --- Scripts SQL, triggers y procedimientos almacenados.
    \item \textbf{Cap\'itulo 4:} Manual de Usuario --- Gu\'ia de uso del sistema.
    \item \textbf{Cap\'itulo 5:} Manual de Despliegue --- Instrucciones de instalaci\'on.
    \item \textbf{Cap\'itulo 6:} Conclusiones --- Objetivos alcanzados y trabajo futuro.
\end{itemize}

